\documentclass[aspectratio=169]{beamer}
\usepackage[english]{babel}
\usepackage[utf8]{inputenc}
\usepackage[T1]{fontenc}
\usepackage{lmodern}
\usepackage{listings}
\usepackage{xcolor}
\usepackage{graphicx}
\usepackage{textcomp}
\DeclareUnicodeCharacter{2014}{---}
\DeclareUnicodeCharacter{2013}{--}
\DeclareUnicodeCharacter{2212}{-}
\DeclareUnicodeCharacter{2192}{\ensuremath{\rightarrow}}
\DeclareUnicodeCharacter{00B1}{\ensuremath{\pm}}
\DeclareUnicodeCharacter{00D7}{\ensuremath{\times}}
\DeclareUnicodeCharacter{00B7}{\ensuremath{\cdot}}
\DeclareUnicodeCharacter{20AC}{EUR}
\DeclareUnicodeCharacter{2070}{\textsuperscript{0}}
\DeclareUnicodeCharacter{00B9}{\textsuperscript{1}}
\DeclareUnicodeCharacter{00B2}{\textsuperscript{2}}
\DeclareUnicodeCharacter{00B3}{\textsuperscript{3}}
\DeclareUnicodeCharacter{2074}{\textsuperscript{4}}
\DeclareUnicodeCharacter{2075}{\textsuperscript{5}}
\DeclareUnicodeCharacter{2076}{\textsuperscript{6}}
\DeclareUnicodeCharacter{2077}{\textsuperscript{7}}
\DeclareUnicodeCharacter{2078}{\textsuperscript{8}}
\DeclareUnicodeCharacter{2079}{\textsuperscript{9}}
\DeclareUnicodeCharacter{2248}{\ensuremath{\approx}}
\DeclareUnicodeCharacter{2264}{\ensuremath{\le}}

% Colors for code
\definecolor{codebg}{RGB}{245,245,245}
\definecolor{codecomment}{RGB}{0,128,0}
\definecolor{codekeyword}{RGB}{0,0,255}
\definecolor{codestring}{RGB}{163,21,21}
\definecolor{bandcolor}{RGB}{200,50,50}

\lstdefinestyle{pythonstyle}{
    language=Python,
    backgroundcolor=\color{codebg},
    basicstyle=\ttfamily\scriptsize,
    keywordstyle=\color{codekeyword},
    commentstyle=\color{codecomment},
    stringstyle=\color{codestring},
    showstringspaces=false,
    breaklines=true,
    frame=single,
    rulecolor=\color{black},
    escapeinside={(*@}{@*)},
    moredelim=**[l][\color{bandcolor}\bfseries]{\#\ ---\ NEW},
    moredelim=**[l][\color{bandcolor}\bfseries]{\#\ ---\ CHANGED},
    literate=
      {á}{{\'a}}1 {é}{{\'e}}1 {í}{{\'i}}1 {ó}{{\'o}}1 {ú}{{\'u}}1
      {Á}{{\'A}}1 {É}{{\'E}}1 {Í}{{\'I}}1 {Ó}{{\'O}}1 {Ú}{{\'U}}1
      {ñ}{{\~n}}1 {Ñ}{{\~N}}1 {ü}{{\"u}}1 {Ü}{{\"U}}1
      {¿}{{?`}}1 {¡}{{!`}}1
      {—}{{---}}1 {–}{{--}}1 {−}{{-}}1
      {±}{{\ensuremath{\pm}}}1 {×}{{\ensuremath{\times}}}1
      {→}{{\ensuremath{\rightarrow}}}1
      {°}{{\textdegree}}1
      {·}{{\ensuremath{\cdot}}}1
      {€}{{\texteuro}}1
      {≈}{{\ensuremath{\approx}}}1
      {≤}{{\ensuremath{\le}}}1
      {⁰}{{\textsuperscript{0}}}1 {¹}{{\textsuperscript{1}}}1
      {²}{{\textsuperscript{2}}}1 {³}{{\textsuperscript{3}}}1
      {⁴}{{\textsuperscript{4}}}1 {⁵}{{\textsuperscript{5}}}1
      {⁶}{{\textsuperscript{6}}}1 {⁷}{{\textsuperscript{7}}}1
      {⁸}{{\textsuperscript{8}}}1 {⁹}{{\textsuperscript{9}}}1
}

\lstset{style=pythonstyle}

% Title slide macro: \lessontitle{Lesson Number}{Title}
\newcommand{\lessontitle}[2]{
    \title{Lesson #1: #2}
    \author{}
    \institute{Tec de Monterrey}
    \begin{frame}[plain]
        \titlepage
    \end{frame}
}

% Figure frame macro: \figureframe{Title}{Image path}
\newcommand{\figureframe}[2]{
    \begin{frame}{#1}
        \begin{center}
            \includegraphics[height=0.8\textheight,width=\textwidth,keepaspectratio]{#2}
        \end{center}
    \end{frame}
}


% Listings pulled from src/ with \lstinputlisting, so the slide is the file.
\lstset{
    basicstyle=\ttfamily\fontsize{8pt}{9.6pt}\selectfont,
    breaklines=true,
    breakatwhitespace=true,
    columns=fullflexible,
    keepspaces=true,
    xleftmargin=0.3em,
    framesep=2pt,
    aboveskip=0pt,
    belowskip=0pt
}

\newcommand{\resultframe}[3]{%
    \begin{frame}{#1}
        \begin{center}
            \includegraphics[height=0.62\textheight,width=\textwidth,keepaspectratio]{#2}
        \end{center}
        \vspace{0.25em}
        {\small #3\par}
    \end{frame}%
}

\newcommand{\claimframe}[2]{%
    \begin{frame}{#1}
        {\small #2\par}
    \end{frame}%
}

\date{}

\begin{document}

\lessontitle{02}{GA Flow}

\section{This lesson}

\begin{frame}{This lesson}
Lesson 01 step 6, with the loop still unnamed, is reproduced here to the digit. Naming the flow must not change the answer.

\vspace{0.6em}
\begin{itemize}
    \item Five named phases; one generation is one function.
    \item The problem is an argument, not welded into the driver.
    \item A cheap stopping rule reports that a run stopped improving --- not that it finished.
\end{itemize}
\end{frame}

% ================= ga_flow_01_the_five_phases =================
\begin{frame}{Mathematical model}
\[
P_{t+1}=R\!\left(M\!\left(C\!\left(S(P_t)\right)\right)\right)
\]
One generation is a random population transition. A patience test such as
\(b_t-b_{t-k}<\varepsilon\) detects little recent improvement; it does not
prove convergence and says nothing about the unknown global optimum.

For \(r\lfloor(T-s|x-x_0|)/h\rfloor\), ordinary one-sided bands have width
\(h/s=0.5\); here the top is \(|x-4|\le0.25\), width 0.5 rather than a point.
\end{frame}

\section{The flow, named}

\begin{frame}{The flow, named}
Lesson 01 finished with a genetic algorithm that worked. Its loop body, though,
was a run of unnamed comment blocks, and you cannot discuss "the flow" when the
flow has no name and no boundary. This step names it.
\end{frame}

\begin{frame}[fragile]{(1) evolve\_one\_generation()}
\lstinputlisting[basicstyle=\ttfamily\fontsize{6pt}{7.2pt}\selectfont,firstline=94,lastline=127]{../src/ga_flow_01_the_five_phases.py}
\end{frame}

\begin{frame}[fragile]{(2) run()}
\lstinputlisting[basicstyle=\ttfamily\fontsize{7pt}{8.5pt}\selectfont,firstline=130,lastline=148]{../src/ga_flow_01_the_five_phases.py}
\end{frame}

\begin{frame}[fragile]{(3) the phase trace}
\lstinputlisting[basicstyle=\ttfamily\fontsize{6pt}{7.2pt}\selectfont,firstline=153,lastline=184]{../src/ga_flow_01_the_five_phases.py}
\end{frame}

\resultframe{The flow, named}{../figures/ga_flow_01_the_five_phases.png}{%
    Same seed, same answer as Lesson 01: \texttt{x=+1.372 f=+0.706}. The run cost \textbf{107} fitness evaluations, not the obvious 110 — an individual selected but neither crossed nor mutated is the same object and is never re-evaluated}

% ================= ga_flow_02_injected_fitness =================
\section{The problem becomes an argument}

\begin{frame}{The problem becomes an argument}
In step 1 the word \texttt{objective} appears inside Individual. That single reference
welds the flow to one problem: to solve a second one you would copy the file.
The book splits \texttt{fitness.py} from \texttt{individual.py} for exactly this reason.
\end{frame}

\begin{frame}[fragile]{(1) the two problems}
\lstinputlisting[basicstyle=\ttfamily\fontsize{8pt}{9.6pt}\selectfont,firstline=47,lastline=61]{../src/ga_flow_02_injected_fitness.py}
\end{frame}

\begin{frame}[fragile]{(2) Individual(genes, f)}
\lstinputlisting[basicstyle=\ttfamily\fontsize{8pt}{9.6pt}\selectfont,firstline=69,lastline=82]{../src/ga_flow_02_injected_fitness.py}
\end{frame}

\begin{frame}[fragile]{(3) the building operators}
\lstinputlisting[basicstyle=\ttfamily\fontsize{8pt}{9.6pt}\selectfont,firstline=99,lastline=115]{../src/ga_flow_02_injected_fitness.py}
\end{frame}

\begin{frame}[fragile]{(4) run(fitness\_function)}
\lstinputlisting[basicstyle=\ttfamily\fontsize{7pt}{8.5pt}\selectfont,firstline=147,lastline=165]{../src/ga_flow_02_injected_fitness.py}
\end{frame}

\begin{frame}[fragile]{(5) the two runs}
\lstinputlisting[basicstyle=\ttfamily\fontsize{7pt}{8.5pt}\selectfont,firstline=168,lastline=190]{../src/ga_flow_02_injected_fitness.py}
\end{frame}

\resultframe{The problem becomes an argument}{../figures/ga_flow_02_injected_fitness.png}{%
    One driver, two problems: the sine run reproduces step 1 exactly, the target run lands \textbf{0.0038} from +4.200. The only difference between the runs is which function was passed in}

% ================= ga_flow_03_population_metrics =================
\section{Measuring the population, not the leader}

\begin{frame}{Measuring the population, not the leader}
Every trace so far reports the best individual. That is the one number a GA is
least honest about: the best can sit still for five generations while the
population underneath it changes completely, or it can be a single lucky outlier
in a population that has learned nothing.
\end{frame}

\begin{frame}[fragile]{(1) population\_metrics()}
\lstinputlisting[basicstyle=\ttfamily\fontsize{8pt}{9.6pt}\selectfont,firstline=67,lastline=83]{../src/ga_flow_03_population_metrics.py}
\end{frame}

\begin{frame}[fragile]{(2) the history table}
\lstinputlisting[basicstyle=\ttfamily\fontsize{6pt}{7.2pt}\selectfont,firstline=188,lastline=230]{../src/ga_flow_03_population_metrics.py}
\end{frame}

\begin{frame}[fragile]{(3) the metrics figure}
\lstinputlisting[basicstyle=\ttfamily\fontsize{7pt}{8.5pt}\selectfont,firstline=232,lastline=257]{../src/ga_flow_03_population_metrics.py}
\end{frame}

\resultframe{Measuring the population, not the leader}{../figures/ga_flow_03_population_metrics.png}{%
    On the sine problem, average fitness climbs from −1.5588 to +0.6228 while gene spread falls from \textbf{6.665 to 0.064} by generation 8. The champion stands still while the population collapses underneath it}

% ================= ga_flow_04_stopping_condition =================
\section{Knowing when to stop}

\begin{frame}{Knowing when to stop}
Every run so far stopped after exactly ten generations, because ten was written
at the top of the file. The book lists three ways to end a run: an acceptable
solution is found, a generation count is reached, or the improvements dry up.
Step 3 measured exactly when the improvements dry up, so the third condition is
now available for free.
\end{frame}

\begin{frame}[fragile]{(1) MAX\_GENERATIONS}
\lstinputlisting[basicstyle=\ttfamily\fontsize{8pt}{9.6pt}\selectfont,firstline=40,lastline=40]{../src/ga_flow_04_stopping_condition.py}
\end{frame}

\begin{frame}[fragile]{(2) PATIENCE, MIN\_IMPROVEMENT}
\lstinputlisting[basicstyle=\ttfamily\fontsize{8pt}{9.6pt}\selectfont,firstline=45,lastline=52]{../src/ga_flow_04_stopping_condition.py}
\end{frame}

\begin{frame}[fragile]{(3) stop on stagnation}
\lstinputlisting[basicstyle=\ttfamily\fontsize{8pt}{9.6pt}\selectfont,firstline=186,lastline=195]{../src/ga_flow_04_stopping_condition.py}
\end{frame}

\begin{frame}[fragile]{(4) the control runs}
\lstinputlisting[basicstyle=\ttfamily\fontsize{6pt}{7.2pt}\selectfont,firstline=273,lastline=326]{../src/ga_flow_04_stopping_condition.py}
\end{frame}

\resultframe{Knowing when to stop}{../figures/ga_flow_04_stopping_condition.png}{%
    The rule saves \textbf{49 to 51} of the 60 allowed generations and costs less than \texttt{MIN\_IMPROVEMENT} on both problems. It also exposes a wart: on 1 of 2 problems the final generation no longer holds the champion — nothing here protects it}

% ================= ga_flow_05_stepped_landscapes =================
\section{Stepped landscapes, and a guess that turned out wrong}

\begin{frame}{Stepped landscapes, and a guess that turned out wrong}
Step 4 priced the stopping rule on two smooth landscapes and found it free. The
natural suspicion is that smoothness was doing the work: on a landscape where
progress arrives in jumps instead of slopes, a flat stretch would mean the
population has not yet found the next step, not that the search is over - and a
patience rule cannot tell those apart.
\end{frame}

\begin{frame}[fragile]{(1) staircase()}
\lstinputlisting[basicstyle=\ttfamily\fontsize{6pt}{7.2pt}\selectfont,firstline=81,lastline=112]{../src/ga_flow_05_stepped_landscapes.py}
\end{frame}

\begin{frame}[fragile]{(2) PROBLEMS}
\lstinputlisting[basicstyle=\ttfamily\fontsize{8pt}{9.6pt}\selectfont,firstline=238,lastline=242]{../src/ga_flow_05_stepped_landscapes.py}
\end{frame}

\begin{frame}[fragile]{(3) the verdict}
\lstinputlisting[basicstyle=\ttfamily\fontsize{6pt}{7.2pt}\selectfont,firstline=365,lastline=436]{../src/ga_flow_05_stepped_landscapes.py}
\end{frame}

\resultframe{Stepped landscapes, and a guess that turned out wrong}{../figures/ga_flow_05_stepped_landscapes.png}{%
    The staircase does \textbf{not} break the rule: the run reaches \texttt{+2.0000}, the landscape's dense-grid reference, and the rule stops it having lost nothing. Every one of the three problems matches its grid reference, shortfall \texttt{+0.0000}}

\section{Next}

\begin{frame}{Next}
On the staircase the patience rule never quits too early: the steps rise, so every boundary is a comparison selection can act on. The hypothesis is refuted, and kept.

\vspace{0.8em}
\textbf{Lesson 03 --- Selection}

The first operator to open, still on one generation of one population. Run-level evidence for pressure is postponed to Lesson 06.
\end{frame}
\end{document}
